\chapter{Trabalhos relacionados}\label{cap:relacionados}

Este capítulo situa o trabalho em três frentes: a aplicação de LLMs ao
domínio de projeto eletrônico, a automação do KiCad por mecanismos
convencionais e o emergente ecossistema de servidores MCP. Ao final,
apresenta-se a lacuna que o trabalho endereça.

\section{LLMs aplicados ao projeto eletrônico}

Liu et al.~\cite{liu2023chipnemo} adaptam modelos de linguagem ao domínio de projeto de
\emph{chips} industriais por meio de tokenização específica, pré-treinamento
continuado e alinhamento com instruções do domínio, avaliando três aplicações:
assistente de engenharia, geração de \emph{scripts} de EDA e análise de
defeitos. Os resultados demonstram que especialização de domínio melhora
tarefas downstream, inclusive superando modelos genéricos em dois dos três
casos avaliados.

Essa linha de trabalho responde à pergunta ``como adaptar o modelo ao
domínio?''. O presente trabalho responde a uma pergunta complementar: ``como
adaptar as ferramentas ao modelo?'', sem alterar o modelo. As duas estratégias
são ortogonais e potencialmente combináveis: um modelo especializado poderia
consumir o mesmo servidor MCP.

\section{Automação do KiCad}

A automação convencional do KiCad apoia-se em três mecanismos
\citep{swig,kicaddev,kicadcli}:

\begin{enumerate}[leftmargin=1.4em,itemsep=0.2em]
  \item \emph{Scripts} Python sobre a API \texttt{pcbnew} (SWIG), típicos para
        geração de geometria, automação de bibliotecas e extração de dados;
  \item \emph{Plugins} executados dentro da interface gráfica, com acesso
        direto ao documento em edição;
  \item Pipelines de \texttt{kicad-cli}, usados em integração contínua para
        verificação e exportação.
\end{enumerate}

Cada abordagem resolve bem o seu caso, mas nenhuma oferece um contrato único
consumível por assistentes: não há, por exemplo, descoberta padronizada de
operações, gestão de sessão entre chamadas ou exposição de estado como
recursos. Uma integração conduzida por IA tende a reproduzir, de forma ad hoc,
esse conjunto de decisões em cada projeto.

\section{O ecossistema de servidores MCP}

Desde a publicação da especificação \citep{mcp2025spec}, servidores MCP foram
criados para domínios variados. No domínio de EDA, o projeto
\texttt{KiCAD-MCP-Server}~\citep{kicadmcpserver}, criado por \emph{mixelpixx}
e aberto a contribuições da comunidade, consolidou uma implementação de
referência em Python e TypeScript. O servidor analisado neste texto é a versão
de trabalho do projeto disponibilizada no \emph{workspace}
\texttt{mcp-kicad-vscode}, que preserva a arquitetura de duas camadas e é
utilizada no estudo de caso conduzido pelo autor
(Capítulo~\ref{cap:estudo-de-caso}).

A literatura acadêmica sobre servidores MCP para EDA ainda é incipiente: os
trabalhos publicados concentram-se no protocolo e em agentes de software, e
não em avaliações de integrações de EDA com métricas de cobertura de
ferramentas, testes e verificação de projeto.

\section{Análise comparativa}

O Quadro~\ref{qua:comparativo} compara as abordagens segundo critérios
relevantes para integração com assistentes de IA.

\begin{quadro}[htbp]
\centering
\small
\caption{Comparação entre abordagens de automação/integração para o fluxo de
PCB no KiCad.}
\label{qua:comparativo}
\begin{tabular}{@{}p{3.4cm}p{3.4cm}p{3.4cm}p{3.4cm}@{}}
\toprule
\textbf{Critério} & \textbf{\emph{Scripts} e plugins} & \textbf{Modelo adaptado ao domínio} & \textbf{Servidor MCP estudado} \\
\midrule
Contrato com o assistente & inexistente (integração ad hoc) & implícito no modelo & padronizado (MCP 2025-06-18) \\
Descoberta de operações & manual (documentação) & embutida no treinamento & \texttt{search\_tools} / categorias \\
Cobertura do fluxo & parcial, por script & parcial, por caso de uso & 233 ferramentas em 24 módulos \\
Verificação & dependente do autor & tarefas avaliadas isoladamente & 2.451 casos de teste; DRC/ERC por \texttt{kicad-cli} \\
Reuso entre assistentes & baixo & baixo (modelo específico) & alto (qualquer cliente MCP) \\
Gratuidade de adaptação & alta (sem protocolo) & alta (treinamento) & média (implementar servidor) \\
\bottomrule
\end{tabular}
\end{quadro}

\section{Lacuna endereçada}

As três frentes analisadas não cobrem, simultaneamente: (i) um contrato
padronizado entre assistentes e o fluxo completo de PCB; (ii) mecanismos de
descoberta e de estado que sustentem uso por agentes; e (iii) evidências
publicadas de qualidade --- testes, cobertura e verificação de regras em um
projeto real. O presente trabalho endereça essa lacuna analisando e
avaliando o \texttt{mcp-kicad-vscode} segundo esse tripé.

\section{Síntese do capítulo}

A literatura demonstra o valor de adaptar modelos ao domínio e de desenhar
interfaces para agentes. Este trabalho explora a segunda direção em um
domínio ainda pouco explorado, com ênfase em amplitude de cobertura,
descoberta, robustez e evidências verificáveis.
